% Annexe de reconstruction numérique, déplacée hors du corps principal.
% =============================================================================
\section{Jeux de paramètres et reconstruction numérique du noyau}
\label{sec:core_parameter_sets}
% =============================================================================

Cette section remplace une dépendance implicite à un fichier de code par une
fiche de reconstruction lisible. Les valeurs intermédiaires sont données dans
la convention du mode local ; les grandeurs de polarisation se déduisent avec
$s=\Omega_0/Z^*$ :
\begin{equation}
\begin{gathered}
A^{(P)}=s^2A^{(u)},\qquad
m_P=s^2M^*,\qquad
\beta^{(P)}=s^4\beta^{(u)},\\
\gamma^{(P)}=s^6\gamma^{(u)},\qquad
\rho^{(P)}=s^8\rho^{(u)},\qquad
G_P=s^{-2}G_u .
\end{gathered}
\label{eq:core_parameter_scaling}
\end{equation}
Pour convertir un coefficient fourni en
$\mathrm{eV/\mathring A^n}$,
\begin{equation}
X\,[\mathrm{Ha/bohr^n}]
=\frac{X\,[\mathrm{eV/\mathring A^n}]}{27.211386245988}
(0.529177210903)^n .
\label{eq:core_ev_angstrom_conversion}
\end{equation}
Tous les chiffres dérivés ci-dessous conservent la précision de leur source ;
des décimales supplémentaires ne constituent pas une information physique
supplémentaire.

\subsection{BaTiO$_3$ : jeu Zhong--Vanderbilt--Rabe}

Le jeu ZVR quartique~\cite{Zhong1995} utilise, en unités atomiques de mode
local,
\begin{equation}
\begin{gathered}
a_0=7.46,\qquad \Omega_0=415.160936,\\
K_2=0.0568,\qquad Z^*=9.956,\qquad
\varepsilon_\infty=5.24,\\
M^*=39.05\ {\rm amu},\\
j_1=-0.02734,\ j_2=0.04020,\ j_3=0.00927,\ j_4=-0.00815,\\
j_5=0.00580,\ j_6=0.00370,\ j_7=0.00185 .
\end{gathered}
\label{eq:core_zvr_input}
\end{equation}
Les entrées élastiques et électrostrictives sont
\begin{equation}
\begin{gathered}
B_{11}=4.64,\quad B_{12}=1.65,\quad B_{44}=1.85,\\
B_{1xx}=-2.18,\quad B_{1yy}=-0.20,\quad B_{4yz}=-0.08,
\end{gathered}
\label{eq:core_zvr_elastic_input}
\end{equation}
les autres composantes étant générées par permutations cubiques. Le potentiel
quartique
$\alpha|\mathbf u|^4+\gamma(u_x^2u_y^2+u_y^2u_z^2+u_z^2u_x^2)$
emploie $\alpha=0.320$ et $\gamma=-0.473$. Après élimination homogène,
\begin{equation}
b_1^{\rm proj}=0.136449,\qquad
b_2^{\rm proj}=0.293874,\qquad
\beta^{\rm proj}_{xxxx}=0.703221 .
\label{eq:core_zvr_quartic_projection}
\end{equation}

La somme d'Ewald et les 26 voisins donnent
\begin{equation}
\begin{array}{c|rrrrr}
&d_0&d_1&d_2&d_3&d_4\\ \hline
\mathrm{ZVR}
&-0.190858109&-0.4180040&-3.2341964&4.48821&0.572574326
\end{array}
\label{eq:core_zvr_dipolar_coeffs}
\end{equation}
et
\begin{equation}
\begin{aligned}
J_0&=0.042200000,\\
c_1&=-3.582850008,\qquad c_2=0.573211480,\qquad
c_3=-1.057380400.
\end{aligned}
\label{eq:core_zvr_sr_derived}
\end{equation}
L'application directe de
Eq.~\eqref{eq:core_five_coefficients} fournit le contrôle final :
\begingroup\small
\begin{equation}
\begin{array}{c|rrrrr}
&A_{0,1}&A_{0,2}&A_{0,3}&A_{0,4}&A_{0,5}\\ \hline
u&-0.035058109&0.155207480&-5.348957200&0.905359992&0.572574326\\
P&-60.96094293&269.88319108&-9301.05712750&1574.28909817&995.62331810
\end{array}
\label{eq:core_zvr_final_kernel}
\end{equation}
\endgroup
avec $s=41.69957172$. Recalculer cette ligne à partir des entrées précédentes
est le premier test unitaire du pipeline.

\subsection{BaTiO$_3$ : jeu WC--GGA Nishimatsu--FERAM}

Le second jeu provient des surfaces d'énergie WC--GGA de
Nishimatsu~\cite{Nishimatsu2010} et de la documentation de génération des
paramètres FERAM~\cite{FERAMParameters}. Les entrées converties sont
\begin{equation}
\begin{gathered}
a_0=7.532372744,\qquad \Omega_0=427.3615143,\\
K_2=0.087822249,\qquad Z^*=10.33,\qquad
\varepsilon_\infty=6.8691464565,\\
M^*=38.24\ {\rm amu},\qquad s=41.37091136,\\
j_1=-0.0214464065,\qquad j_2=-0.0116188030,\\
j_3=0.0070951144,\qquad j_4=-0.0062912217,\\
j_5=0,\qquad j_6=0.0028495045,\qquad j_7=0 .
\end{gathered}
\label{eq:core_nish_input}
\end{equation}
Les coefficients élastiques sources, avant conversion, sont
\begin{equation}
\begin{aligned}
B_{11}&=126.7316715\ {\rm eV},&
B_{12}&=41.7582964\ {\rm eV},\\
B_{44}&=49.2408864\ {\rm eV},&
B_{1xx}&=-185.3471876\ {\rm eV/\mathring A^2},\\
B_{1yy}&=-3.2809295\ {\rm eV/\mathring A^2},&
B_{4yz}&=-14.5501739\ {\rm eV/\mathring A^2}.
\end{aligned}
\label{eq:core_nish_elastic_source}
\end{equation}
Le potentiel local emploie
\begin{equation}
\begin{gathered}
\alpha=0.2276183984,\qquad \gamma=-0.3327945785,\\
k_1=-0.2162513067,\qquad k_2=0.1593766928,\\
k_3=0.6699445355,\qquad k_4=0.1450680738,
\end{gathered}
\label{eq:core_nish_anharmonic}
\end{equation}
en $\mathrm{Ha/bohr^4}$, $\mathrm{Ha/bohr^6}$ et
$\mathrm{Ha/bohr^8}$ selon l'ordre. Après projection homogène,
$b_1^{\rm proj}=0.1119078117$ et $b_2^{\rm proj}=0.1158966279$.

Les coefficients dérivés sont
\begin{equation}
\begin{aligned}
d_0&=-0.1522618376,&d_1&=-0.3399749261,&
d_2&=-2.6304665897,\\
d_3&=3.6503913681,&d_4&=0.4567855127,\\
J_0&=-0.0546311676,&c_1&=-2.0765766351,\\
c_2&=0.4788914499,&c_3&=0.
\end{aligned}
\label{eq:core_nish_derived}
\end{equation}
On obtient
\begingroup\small
\begin{equation}
\begin{array}{c|rrrrr}
&A_{0,1}&A_{0,2}&A_{0,3}&A_{0,4}&A_{0,5}\\ \hline
u&-0.03124850714&0.1389165238&-2.6304665897&1.5738147330&0.4567855127\\
P&-53.48345447&237.76289673&-4502.18115849&2693.66623606&781.81229778
\end{array}
\label{eq:core_nish_final_kernel}
\end{equation}
\endgroup
À $500\,\mathrm K$ et pression nulle, la fermeture locale projetée
simplifiée, sur grille $48^3$, donne
\begin{equation}
\begin{aligned}
A_1^{(u)}&=0.002930749645,&
\Delta A_1^{(4)}&=0.055301815578,\\
\Delta A_1^{(6)}&=-0.028554029685,&
\Delta A_1^{(8)}&=0.007431470895.
\end{aligned}
\label{eq:core_nish_500k_check}
\end{equation}
Ce contrôle concerne le solveur \texttt{scha\_paraelc.py}, qui traite
$\beta_{\rm proj}=\beta-\Lambda^{\rm hom}/2$ comme un vertex local.
Il ne valide pas la fermeture finale à secteurs homogène et inhomogène
séparés : celle-ci conserve le terme homogène collectif
$-\Lambda^{\rm hom}G/2$, au lieu de la contribution locale projetée
$-3\Lambda^{\rm hom}G/2$, ainsi que la convolution acoustique distincte.
Le passage de $40^3$ à $48^3$ déplace $A_1^{(u)}$ de
$0.0029264173$ à $0.0029307496$.

\subsection{SrTiO$_3$ : jeu Nishimatsu 2016}

Le tableau complet pour SrTiO$_3$ a été publié ultérieurement
~\cite{Nishimatsu2016}. Les entrées de structure et de dynamique sont
\begin{equation}
\begin{gathered}
a_0=3.901\ \mathring{\rm A},\quad
M^*_{\rm dip}=43.61\ {\rm amu},\quad
M^*_{\rm ac}=36.70\ {\rm amu},\\
Z^*=9.28,\qquad \varepsilon_\infty=6.46,\qquad p=0.
\end{gathered}
\label{eq:core_sto_basic_input}
\end{equation}
Les coefficients harmoniques sources, en
$\mathrm{eV/\mathring A^2}$, sont
\begin{equation}
K_2=10.316,\quad
(j_1,\ldots,j_7)=(-2.012,-1.815,0.590,-0.567,0,0.238,0).
\label{eq:core_sto_harmonic_source}
\end{equation}
Les valeurs élastiques sont
\begin{equation}
\begin{aligned}
B_{11}&=131.33\ {\rm eV},&
B_{12}&=36.26\ {\rm eV},&
B_{44}&=41.30\ {\rm eV},\\
B_{1xx}&=-102.09\ {\rm eV/\mathring A^2},&
B_{1yy}&=0.5299\ {\rm eV/\mathring A^2},&
B_{4yz}&=-15.494\ {\rm eV/\mathring A^2}.
\end{aligned}
\label{eq:core_sto_elastic_source}
\end{equation}
La source arrondit fortement ces nombres. Le contact reconstruit
$2K_2+d_0+J_0=-0.0024173835\,\mathrm{Ha/bohr^2}$ diffère ainsi de la courbure
indépendamment tabulée
$2\kappa=-0.0025932963\,\mathrm{Ha/bohr^2}$. Cette dernière est retenue pour la
masse à $\Gamma$, sans modifier la dispersion périodique hors de $\Gamma$.
Avec $s=43.16944605$, les coefficients de diagnostic à petit $k$ sont
\begingroup\small
\begin{equation}
\begin{array}{c|rrrrr}
&A_{0,1}&A_{0,2}&A_{0,3}&A_{0,4}&A_{0,5}\\ \hline
u&-0.0025932963&0.2689681617&-2.3065100079&1.7965640186&0.4181661231\\
P&-4.8328697591&501.24935461&-4298.41452401&3348.0786316&779.29483544
\end{array}
\label{eq:core_sto_final_kernel}
\end{equation}
\endgroup
Cette incohérence de précision est conservée comme incertitude de source et
n'est pas absorbée silencieusement dans une pression ajustable.

\subsection{Ordre de reconstruction à partir des tableaux}

Un lecteur qui repart de ces fiches doit obtenir les mêmes objets dans l'ordre
suivant :
\begin{enumerate}[label=\textbf{R\arabic*.},leftmargin=3.2em,itemsep=2pt]
  \item convertir $a_0$, $K_2$, $j_i$, $B_{\ell m}$,
  $B_{\ell\alpha\beta}$ et les coefficients anharmoniques ;
  \item calculer $\Omega_0$, $s$, $m_P$ et
  $\Lambda^{\rm hom}=B^T\mathcal B^{-1}B$ ;
  \item générer les 26 matrices $J(\mathbf n)$, puis
  $\Phi^{\rm SR}(\mathbf k)$ ;
  \item calculer la somme d'Ewald et vérifier $d_4$ par
  $4\pi Z^{*2}/(\varepsilon_\infty\Omega_0)$ ;
  \item assembler $\Phi^{\rm harm}$ et retrouver les cinq coefficients des
  Eqs.~\eqref{eq:core_zvr_final_kernel},
  \eqref{eq:core_nish_final_kernel} ou
  \eqref{eq:core_sto_final_kernel} dans la limite $k\to0$ ;
  \item seulement alors lancer la fermeture de
  Eq.~\eqref{eq:core_scalar_residual}.
\end{enumerate}
